\documentclass[11pt]{article}
\usepackage[margin=1in]{geometry}
\usepackage{amsmath,amssymb}
\usepackage[hidelinks]{hyperref}
\usepackage{microtype}

\newcommand{\T}{\mathbb T}
\newcommand{\spectrum}{\sigma}

\title{A Proper-Arc Quasianalytic Contraction\\
\large The question in arXiv:1712.07075 is answered by arXiv:2005.04031}
\author{Literature-resolution packet for \texttt{fa\_banach\_001}}
\date{17 August 2026}

\begin{document}
\maketitle

\noindent\textbf{Status.}
Literature already answered.  This packet matches the exact quantifiers of
the 2017 question to a 2020 construction; it does not claim a new proof.

\section{The exact question}

For a quasianalytic contraction $T$, write $\pi(T)$ for its quasianalytic
spectral set and $\spectrum(T)$ for its spectrum.  Gamal' records the
following question in the introduction of the source paper
\cite[PDF p.~3]{source}:

\begin{quote}
Does there exist a contraction $T$ such that
\[
                 \pi(T)=\spectrum(T)\ne\T ?
\]
\end{quote}

The surrounding discussion explains the gap.  Earlier examples with
$\pi(T)=\spectrum(T)\cap\T\ne\T$ had spectrum extending off the unit circle,
whereas the source's new example had $\spectrum(T)=\T$ and
$\pi(T)\ne\T$.  Neither supplied equality on a proper subset of $\T$.

\section{The later answer}

The same author later constructed exactly the missing example in
\emph{Example of quasianalytic contraction whose spectrum is a proper subarc
of the unit circle}, arXiv:2005.04031 \cite{answer}.  Its Corollary~9.7 states
\cite[Corollary~9.7, PDF p.~27]{answer}:

\begin{quote}
There exists a quasianalytic contraction $R$ with
\[
 \spectrum(R)=\{e^{it}:t\in[0,\pi]\}
\]
and with unitary asymptote $U_{\spectrum(R)}$.  Consequently,
$\spectrum(R)$ coincides with the quasianalytic spectral set of $R$ and
$\spectrum(R)\ne\T$.
\end{quote}

Therefore $R$ is a contraction and
\[
 \boxed{\ \pi(R)=\spectrum(R)=\{e^{it}:0\le t\le\pi\}\ne\T\ },
\]
which is a positive answer to every clause of the exact 2017 existence
question.

The proof of the later corollary uses the operator constructed in its
Theorem~9.6.  It produces a square root $R$ of an operator $T$ with
$\spectrum(T)=\T$, controls $\spectrum(R)$ to lie in the closed upper
half-plane, and uses spectral mapping plus the unitary asymptote to force the
closed semicircle.  Its final similarity step gives the stated quasianalytic
contraction.

\section{\hspace{0.35em}Why the later paper says ``partial answer''}

The abstract of arXiv:2005.04031 calls the construction a partial answer to
the broader K\'erchy--Szalai Question~2 because it gives no estimate of
$\lVert R^{-1}\rVert$ \cite{answer}.  This is an additional quantitative
refinement.  The 2017 source sentence asks only for existence of a
contraction with $\pi(T)=\spectrum(T)\ne\T$, and Corollary~9.7 settles that
sentence completely.

The source also separately says that the author did not know whether its own
2017 operator satisfies the hypotheses of one of its hyperinvariant-subspace
criteria.  The 2020 proper-arc construction resolves the flagged spectral
existence question, not that distinct operator-specific check.

\section{\hspace{0.35em}Verification and disposition}

The exact source sentence was checked in the stored source and on PDF page 3.
The later abstract, introduction, Theorem~9.6, and Corollary~9.7 were checked
in the stored arXiv source and in the downloaded PDF.  The answer is direct,
not an inference from an adjacent theorem.  The target should therefore be
archived as literature already answered rather than used for a new-proof
attempt.

\begin{thebibliography}{9}
\bibitem{source}
M. F. Gamal',
\emph{Some sufficient conditions for existence of hyperinvariant subspaces
for operators intertwined with unitaries},
arXiv:1712.07075, 2017.

\bibitem{answer}
M. F. Gamal',
\emph{Example of quasianalytic contraction whose spectrum is a proper subarc
of the unit circle},
arXiv:2005.04031, 2020.
\end{thebibliography}

\end{document}
